\documentclass{article}
\usepackage{graphicx}
\usepackage{amsmath}
\usepackage{amsfonts}
\usepackage{bookmark}
\usepackage{booktabs}
\usepackage{exercises}
\graphicspath{ {./images/} }
\usepackage{indentfirst}
\usepackage{amsthm}
\usepackage{yhmath}
\usepackage{placeins}
\newcommand*{\oldneg}{\mathord{\sim}}
\newcommand\showdiv[1]{\overline{\smash{)}#1}}
\newcommand\Mydiv[1]{%
\smash{\raise.3ex\hbox{$\big)$}}\mkern-8mu
        \overline{\enspace\strut#1}}
\setlength{\tabcolsep}{5pt}
\theoremstyle{definition}
\newtheorem{definition}{Definition}

\setlength{\parindent}{.5in}
\setlength{\parskip}{12pt}
\renewcommand{\arraystretch}{2}
\newcommand{\?}{\stackrel{?}{=}}

\title{4.6 Answers}
\begin{document}
\maketitle

\section*{1}
The equation $5^x=8$ is an example of an exponential equation.

\section*{2}
Logarithmic equations may have extraneous solutions when one of the solutions is outside of the domain of the logarithmic function. For example, 
a negative solution for the equation $\log x=a$ would not be possible, as $\log x$ where $x\le 0$ is not defined.

\section*{8}
\[ 3^{3x-7}=81^{12-3x} \]
\[ 3^{3x-7}=(3^4)^{12-3x} \]
\[ 3^{3x-7}=3^{4(12-3x)} \]
\[ 3x-7=4(12-3x) \]
\[ 3x-7=48-12x \] \[ 15x=55 \] \[ x=\frac{11}{3} \]

\section*{9}
\[ 36^{5x+2}=(\frac{1}{6})^{11-x} \] \[ (6^2)^{5x+2}=(6^{-1})^{11-x} \] \[ 6^{10x+4}=6^{x-11} \] \[ 10x+4=x-11 \] \[ x=-\frac{5}{3} \]

\section*{10}
\[ 10^{3x-10}=(\frac{1}{100})^{6x-1} \] \[ 10^{3x-10}=(10^{-1})^{6x-1} \] \[ 10^{3x-10}=10^{1-6x} \] \[ 3x-10=1-6x \] \[ x=\frac{11}{9} \]

\section*{11}
\[ 25^{10x+8}=(\frac{1}{125})^{4-2x} \] \[ (5^2)^{10x+8}=(5^{-3})^{4-2x} \] \[ 5^{20x+16}=5^{6x-12} \] \[ 20x+16=6x-12 \] \[ x=-2 \]

\section*{20}
\[ 0.5^x-0.25=4 \] \[ \log_{0.5}(0.5^x)=\log_{0.5} 4.25 \] \[ x=\dfrac{\log{4.25}}{\log{0.5}}\approx -2.087 \]

\section*{21}
\[ \frac{1}{3}(6)^{-4x}+1=6 \] \[ 6^{-4x}=15 \] \[ \log_6 (6^{-4x})=\log_6 15 \] \[ -4x=\dfrac{\log{15}}{\log{6}} \] \[ -4x\approx 1.51 \] 
\[ x\approx -6.05 \]

\section*{22}
\[ 2^{0.1x}-5=7 \] \[ \log_2 (2^{0.1x})=\log_2 12 \] \[ 0.1x=\dfrac{\log 12}{\log 2} \] \[ x\approx 35.8 \]

\section*{23}
\[ \frac{3}{4}e^{2x}+\frac{7}{2}=4 \] \[ e^{2x}=\frac{2}{3} \] \[ \ln (e^{2x})=\ln\frac{2}{3} \] \[ 2x=\frac{\log\frac{2}{3}}{\log e} \]
\[ x\approx -0.2 \]

\section*{28}
\[ \log\left( 12x-11\right) =\log\left( 3x+13\right) \] \[ 12x-11=3x+13 \] \[ x=\frac{8}{3} \]
Check for extraneous solutions:
\[ \log\left( 12(\frac{8}{3}-11\right)\?\left( 3\frac{8}{3}+13\right) \] \[ \log\left( 13\right)\ne\log\left( 21\right) \]

\section*{31}
\[ \log_8\left( 5-12x\right)=\log_8\left(6x-1\right) \] \[ 5-12x=6x-1 \] \[ x=\frac{1}{3} \]
Check for extraneous solutions:
\[ \log_8\left( 5-12(\frac{1}{3}\right)=\log_8 1=0 \] and \[ \log_8\left(6\frac{1}{3}-1\right)=1 \] meaning $x=\frac{1}{3}$ really is a solution.

\section*{41}
\[ \log_5\left( x+4\right)+\log_5\left(x+1\right)=2 \] \[ \log_5\left(x^2+5x+4\right)=2 \]
\[ 5^{\log_5\left(x^2+5x+4\right)}=5^2 \] \[ x^2+5x+4=25 \] \[ x^2+5x+(\frac{5}{2})^2=31.25 \] \[ (x+\frac{5}{2})^2=31.25 \]  
\[ x=-\frac{5}{2}\pm\sqrt{31.25}=53.422\vee -58.422 \]
Check for extraneous solutions:

\section*{43}
\[ \log_3(x-9)+\log_3(x-3)=2 \] \[ \log_3(x^2-12x+27)=2 \] \[ 3^{\log_3(x^2-12x+27)}=3^2 \] \[ x^2-12x+27=9 \]

\section*{44}
A

\section*{48}
\[ 3^{x+4}=6^{2x-5} \] \[ 3^{x+4}=3^{4x-10} \] \[ x+4=4x-10 \] \[ x=\frac{14}{3} \]

\section*{49}
\[ 10^{3x-8}=2^{5-x} \] \[ \log_10 (10^{3x-8})=\log_10 (2^{5-x}) \] \[ 3x-8=\frac{\log_2 (2^{5-x})}{\log_2 10} \]
\[ \log_2 10\times 3x-8=5-x \] \[ \frac{\log 10}{\log 2}\times 3x-8=5-x \] \[ 9.96x-26.56=5-x \] \[ x\approx 3.169 \]

\section*{50}
\[ \log_2 (x+1)=\log_8 3x \] \[ \log_2 (x+1)=\frac{\log_2 3x}{\log_2 8} \] \[ 3\log_2 (x+1)=\log_2 3x \] \[ \log_2 (x+1)^2=\log_2 3x \]
\[ (x+1)^2=3x \] \[ x^2+2x+1=3x \] \[ \]

\section*{51}
\[ \log_3 x=\log_9 6x \] \[ \log_3 x=\frac{\log_3 6x}{\log_3 9} \] \[ \log_3 x=\frac{\log_3 6x}{2} \] \[ 2\log_3 x=\log_3 6x \]
\[ \log_3 x^2=\log_3 6x \] \[ x^2=6x \] \[ x=6 \]

\section*{52}
\[ 2^{2x}-12\times 2^x+32=0 \] \[ 2^{2x}-8\times 2^x-4\times 2^x+32=0 \] \[ 2^x(2^x-8)-4(2^x-8)=0 \] \[ (2^x-4)(2^x-8)=0 \]
\[ x=\log_2 4=2\vee x=\log_2 8=3 \]

\section*{53}
\[ 5^{2x}+20\times 5^x-125=0 \] \[ 5^{2x}+20\times 5^x=125 \] \[ 5^{2x}+20\times 5^x+(\frac{20}{2})^2=225 \]
\[ (5^x+10)^2=225 \] \[ 5^x+10=15 \] \[ 5^x=5 \] \[ \log_5 (5^x)=\log_5 5 \] \[ x=5 \]

\end{document}
